Todo el experimento está en el repositorio \url{\Repo}, etiqueta \texttt{\Tag}, que es la versión evaluada en este informe. El \texttt{README} basta para repetirlo en otra máquina.

\textbf{Entorno.} SUMO 1.12 en la imagen Docker \texttt{van3twin:jammy} del repositorio del curso; Python 3.13.12 con numpy 2.5.4, pandas 3.1.1, pyarrow 26.0.1, scikit-learn 1.9.1, PyTorch 2.14.1 (CPU), joblib 1.6.0, matplotlib 3.11.2 y pytest 9.2.0, fijados en \texttt{requirements.txt}. El informe se compila con \texttt{tectonic}.

\textbf{Semillas y parámetros.}
\begin{itemize}
\item Calendario de días: semilla 2026. Cada día en SUMO y \texttt{duarouter}: semilla 1000~+~día. Días extra del AID: semilla 2027. Barrido de saturación: semillas 101, 202 y 303.
\item GCN-GRU: \texttt{torch.manual\_seed(42)}, AdamW con tasa 0,002, 32 unidades ocultas, lotes de 16, 12 épocas con sensores y 8 con V2X.
\item AID: \texttt{random\_state=42}, 300 iteraciones, tasa 0,06 y 48 hojas.
\item SIL: radio de 2~km, penalización de 600~s, caducidad de 3~min y calentamiento de 60~min.
\end{itemize}
El Anexo~\ref{an:param} lista los parámetros del generador.

\textbf{Un solo script.} \texttt{./repetir\_experimento.sh} entrena y evalúa todo desde el dataset (unas 2,5~h en CPU). Con \texttt{--desde-cero} construye la red desde OSM y genera el dataset en SUMO (unas 13~h). Con \texttt{--solo-figuras} regenera en 2~min las tablas, figuras, el tablero y las pruebas desde las salidas guardadas en \texttt{v2/resultados/}. La Tabla~\ref{tab:scripts} indica qué genera cada resultado.

\begin{table}[h]
\centering
\caption{Scripts que regeneran cada tabla y figura}
\label{tab:scripts}
\scriptsize
\setlength{\tabcolsep}{2.5pt}
\begin{tabular}{@{}L{2.35cm}L{2.9cm}L{3.1cm}@{}}
\toprule
Resultado & Script & Salida \\
\midrule
Tabla~\ref{tab:pred}, Fig.~\ref{fig:horizonte} & \texttt{entrenar\_prediccion.py} & \texttt{metricas\_prediccion\_*.json} \\
Tabla~\ref{tab:aid}, Fig.~\ref{fig:aidazar} & \texttt{entrenar\_aid.py} & \texttt{metricas\_aid\_*.json} \\
Tabla~\ref{tab:sil}, Fig.~\ref{fig:silcola} & \texttt{sil\_sumo.py} & \texttt{sil\_inc119\_3\_gt.json} \\
Tabla~\ref{tab:valida} & \texttt{validar\_sintetico.py} & \texttt{validacion\_sintetico.json} \\
Figuras (PDF) & \texttt{informe/figuras\_informe.py} & \texttt{informe/figuras/} \\
Fig.~\ref{fig:dash} & \texttt{construir\_dashboard.py} & \texttt{dashboard/} \\
Conteo de pruebas & \texttt{pytest v2/pruebas} & \texttt{informe/pruebas.tex} \\
\bottomrule
\end{tabular}
\end{table}

\textbf{Datos.} El dataset completo (1,5~GB) se regenera de forma determinista con el script. El repositorio incluye la red, el OSM de origen, las tablas pequeñas del dataset y todas las salidas que respaldan las cifras.
